\documentclass[10pt,UTF8]{article}

\usepackage{geometry}

\usepackage{ctex}

\geometry{a4paper,scale=0.9}
\ctexset{today=old}

\usepackage[bookmarksnumbered=true]{hyperref}  % 生成大纲


\usepackage{times}  % 设置正文字体为 Adobe Times Roman

% \usepackage{makecell}
% \usepackage{multirow} % 表格多行合并

% \usepackage{graphicx} %插入图片的宏包
% \usepackage{subfigure}
% \usepackage{float} %设置图片浮动位置的宏包

\usepackage{amsmath}
\usepackage{amssymb}
\usepackage{mathtools}

\usepackage{physics}

% \usepackage{siunitx}
% % 关闭警告
% \ExplSyntaxOn
% \msg_redirect_name:nnn { siunitx } { physics-pkg } { none }
% \ExplSyntaxOff

%%%%%%%%%%%%%%%%%%%%%%%% tcolorbox %%%%%%%%%%%%%%%%%%%%%%%%%%
\usepackage[most]{tcolorbox}

\newenvironment{definition box}
{\begin{tcolorbox}[colback=blue!5!white,colframe=blue!75!black,parbox=false,before upper=\par,before lower=\par]}
{\end{tcolorbox}}

\newenvironment{theorem box}
{\begin{tcolorbox}[colback=yellow!5!white,colframe=yellow!75!black,parbox=false,before upper=\par,before lower=\par]}
{\end{tcolorbox}}
%%%%%%%%%%%%%%%%%%%%%%%%%%%%%%%%%%%%%%%%%%%%%%%%%%%%%%%%%%%%
            
%%%%%%%%%%%%%%%%%%%%% amsthm %%%%%%%%%%%%%%%%
\usepackage{amsthm}

\theoremstyle{definition}
\newtheorem{definition inner}{Definition}[section]
\newenvironment{definition}[1][]
{\begin{definition box}\begin{definition inner}[#1]\hfill\par}
{\end{definition inner}\end{definition box}}

\theoremstyle{plain}
\newtheorem{theorem inner}{Theorem}[section]
\newenvironment{theorem}[1][]
{\begin{theorem box}\begin{theorem inner}[#1]\hfill\par}
{\end{theorem inner}\end{theorem box}}

\theoremstyle{definition}
\newtheorem*{proof inner}{\textit{Proof}}
\renewenvironment{proof}[1][]
{\begin{proof inner}[#1]\hfill\par}
{\qed\end{proof inner}}

\theoremstyle{remark}
\newtheorem*{remark}{\textit{Remark}}
%%%%%%%%%%%%%%%%%%%%%%%%%%%%%%%%%%%%%%%%%%%%%%


\newcommand{\mycases}[2][0.8]{
    \left\{\vphantom{\scalebox{#1}{
        $\begin{matrix}#2\end{matrix}$
    }}\right.\!\!\!\!
    \begin{array}{l@{\quad}l}#2\end{array}
}

\newcommand{\ju}{\mathrm{j}}  % 虚数单位
\newcommand{\e}{\mathrm{e}{\mkern 1 mu}}

\newcommand{\Ev}{\,\mathcal{E}\ensuremath{v}\qty}
\newcommand{\Od}{\,\mathcal{O}\ensuremath{d}\qty}

\renewcommand{\Re}{\,\mathcal{R}\ensuremath{e}\qty}
\renewcommand{\Im}{\,\mathcal{I}\ensuremath{m}\qty}

\newcommand{\sinc}{\operatorname{sinc}} 
\newcommand{\rect}{{\mkern 1.5 mu}\operatorname{rect}\!\qty} % 矩形函数



\begin{document}

\part*{The Continuous-Time Fourier Transform}

\section{From Periodic to Aperiodic Signal}

Consider an aperiodic signal $x(t)$
that is nonzero only over $[0,T_1]$.
There is a corresponding signal $\tilde{x}(t)$ with period $T$
that is the same as $x(t)$ on $[0,T_1]$ and $T\geqslant T_1$.
\begin{equation*}
    x(t) = \mycases{
        x(t) \qfor t\in[0,T_1] \\
        0 \qotherwise
    }
    \qquad\qquad
    \tilde{x}(t+k T) = x(t)
    \qfor k\in\mathbb{Z}\qand t\in[0,T]
    \qquad\qquad
    T_1 \leqslant T
\end{equation*}
Intuitively, we have
\begin{equation*}
    x(t) = \lim_{T\to\infty} \tilde{x}(t)
\end{equation*}

The Fourier coefficients for the periodic signal $\tilde{x}(t)$ are
\begin{equation*}
\tilde{a}_k = 
\frac{1}{T}\int_{0}^{T}\tilde{x}(t)\e^{-\ju k\omega_0 t}\dd{t}
=\frac{1}{T}\int_{0}^{T}x(t)\e^{-\ju k\omega_0 t}\dd{t}
=\frac{1}{T}\int_{-\infty}^{\infty}x(t)\e^{-\ju k\omega_0 t}\dd{t}
\end{equation*}
We define
\begin{equation*}
    X(\ju\omega) :=
    \int_{-\infty}^{\infty}x(t)\e^{-\ju\omega t}\dd{t}
    \qquad\qquad
    X(\ju k\omega_0) =
    \int_{-\infty}^{\infty}x(t)\e^{-\ju k\omega_0 t}\dd{t}
    = T\tilde{a}_k
\end{equation*}
Then
\begin{equation*}
    \tilde{x}(t) = 
    \sum_{k=-\infty}^{\infty}
    \tilde{a}_k \e^{\ju k\omega_0 t}
    = \sum_{k=-\infty}^{\infty}
    \frac{1}{T}X(\ju k\omega_0)\e^{\ju k\omega_0 t}
    = \sum_{k=-\infty}^{\infty}\frac{\omega_0}{2\pi}
    X(\ju k\omega_0)\e^{\ju k\omega_0 t}
\end{equation*}
\begin{align*}
    x(t) &= \lim_{T\to\infty} \tilde{x}(t) \\
    &= \lim_{\omega_0\to 0^+} \frac{1}{2\pi}\sum_{k=-\infty}^{\infty}
    X(\ju k\omega_0)\e^{\ju k\omega_0 t}\omega_0 \\
    &= \frac{1}{2\pi}\int_{-\infty}^{\infty}
    X(\ju\omega)\e^{\ju\omega t}\dd{\omega}
\end{align*}

\begin{definition}[The Fourier Transform Pair of a CT Aperiodic Signal]
If an aperiodic signal $x(t)$ has a Fourier transform, then
\begin{align*}
    x(t) &=
    \frac{1}{2\pi}\int_{-\infty}^{\infty}
    X(\ju\omega)\e^{\ju\omega t}\dd{\omega}
    \tag{The Synthesis Equation} \\
    X(\ju\omega) &=
    \int_{-\infty}^{\infty}x(t)\e^{-\ju\omega t}\dd{t}
    \tag{The Analysis Equation}
\end{align*}
where the function $X(\ju\omega)$ is called
the \emph{Fourier transform} or \emph{Fourier integral} of $x(t)$,
and the \emph{synthesis equation} is called the \emph{inverse Fourier transform}.
\end{definition}

\begin{theorem box}
The Fourier coefficients of the periodic signal $\tilde{x}(t)$
are proportional to spaced samples of
the Fourier transform of the finite-duration signal $x(t)$:
\begin{equation*}
    \tilde{a}_k =
    \left.\frac{1}{T}X(\ju\omega)\right|_{\omega=k\omega_0}
\end{equation*}
\end{theorem box}
For aperiodic signals, $X(\ju\omega)\dfrac{\dd{\omega}}{2\pi}$
can be seen as ``amplitude'' of the complex exponential compoents.
Therefore, the transform $X(\ju\omega)$ is referred to as
the \emph{spectrum} of $x(t)$.



\section{Fourier Transform for Periodic Signals}

\begin{theorem}
    For a periodic signal $x(t)$ whose FS coefficients are $\{a_k\}$,
    the Fourier transform $X(\ju\omega)$ is given by
    \begin{align*}
        X(\ju\omega) &= \sum_{k=-\infty}^{\infty} 
        2\pi a_k \delta(\omega-k\omega_0) \\
        &= \mycases{
            2\pi a_k\delta(0) &\qfor \frac{\omega}{\omega_0}=k\in\mathbb{Z} \\
            0 &\qotherwise
        }
    \end{align*}
\end{theorem}
\begin{proof}
\begin{align*}
    X(\ju\omega) 
    \quad\xrightarrow{\mathcal{F}^{-1}}\quad
    \frac{1}{2\pi}\int_{-\infty}^{\infty}
    X(\ju\omega)\e^{\ju\omega t}\dd{\omega}
    &=  \frac{1}{2\pi}\int_{-\infty}^{\infty}
    \sum_{k=-\infty}^{\infty}
    2\pi a_k \delta(\omega-k\omega_0)
    \e^{\ju\omega t}\dd{\omega} \\
    &= \sum_{k=-\infty}^{\infty}
    a_k \int_{-\infty}^{\infty}
    \delta(\omega-k\omega_0)
    \e^{\ju\omega t}\dd{\omega} \\
    &= \sum_{k=-\infty}^{\infty}
    a_k\e^{\ju k\omega_0 t} \\
    &= x(t)
\end{align*}
\end{proof}

\section{Properties of CT Fourier Transform}

\subsection{}
Suppose we have:
\begin{gather*}
    \mycases[0.7]{
        x(t) \xrightleftharpoons[\mathcal{F}^{-1}]{\mathcal{F}}
        X(\ju\omega) \\
        y(t) \xrightleftharpoons[\mathcal{F}^{-1}]{\mathcal{F}}
        Y(\ju\omega)
    }
\end{gather*}
Then we have the following results:
\begin{align*}
\shortintertext{Linearity}
    a\cdot x(t) + b\cdot y(t)
    \quad&\xrightleftharpoons[\mathcal{F}^{-1}]{\mathcal{F}}\quad
    a\cdot X(\ju\omega) + b\cdot Y(\ju\omega)
\shortintertext{Time Shifting}
    x(t-t_0)
    \quad&\xrightleftharpoons[\mathcal{F}^{-1}]{\mathcal{F}}\quad
    \e^{-\ju\omega t_0} X(\ju\omega)
\shortintertext{Frequency Shifting}
    \e^{\ju\omega_0 t} x(t)
    \quad&\xrightleftharpoons[\mathcal{F}^{-1}]{\mathcal{F}}\quad
    X(\ju(\omega-\omega_0))
\shortintertext{Time Reversal}
    x(-t)
    \quad&\xrightleftharpoons[\mathcal{F}^{-1}]{\mathcal{F}}\quad
    X(-\ju\omega)
\shortintertext{Scaling}
    x(a t)
    \quad&\xrightleftharpoons[\mathcal{F}^{-1}]{\mathcal{F}}\quad
    \frac{1}{\abs{a}}X(\frac{\ju\omega}{a})
\shortintertext{Conjugation}
    x^\ast(t)
    \quad&\xrightleftharpoons[\mathcal{F}^{-1}]{\mathcal{F}}\quad
    X^\ast(-\ju\omega)
\shortintertext{Convolution}
    x(t)\ast y(t)
    \quad&\xrightleftharpoons[\mathcal{F}^{-1}]{\mathcal{F}}\quad
    X(\ju\omega) Y(\ju\omega)
\shortintertext{Multiplication}
    x(t) y(t)
    \quad&\xrightleftharpoons[\mathcal{F}^{-1}]{\mathcal{F}}\quad
    \frac{1}{2\pi}X(\ju\omega)\ast Y(\ju\omega) =
    \frac{1}{2\pi}\int_{-\infty}^{\infty}
    X(\ju\theta)Y(\ju\omega-\ju\theta)\dd{\theta}
\shortintertext{Differentiation in Time}
    \dv{x(t)}{t}
    \quad&\xrightleftharpoons[\mathcal{F}^{-1}]{\mathcal{F}}\quad
    \ju\omega X(\ju\omega)
\shortintertext{Integration in Time}
    \int_{-\infty}^t x(\tau)\dd{\tau} = x(t) \ast u(t)
    \quad&\xrightleftharpoons[\mathcal{F}^{-1}]{\mathcal{F}}\quad
    \frac{1}{\ju\omega} X(\ju\omega) + \pi X(0)\delta(\omega)
\shortintertext{Differentiation in Frequency}
    t\cdot x(t)
    \quad&\xrightleftharpoons[\mathcal{F}^{-1}]{\mathcal{F}}\quad
    \ju\dv{\omega} X(\ju\omega)
\shortintertext{Even-Odd Decomposition for Real Signals}
    x(t)\qq*{ real:} \Ev{x(t)}
    \quad&\xrightleftharpoons[\mathcal{F}^{-1}]{\mathcal{F}}\quad
    \Re{X(\ju\omega)} \\
    x(t)\qq*{ real:} \Od{x(t)}
    \quad&\xrightleftharpoons[\mathcal{F}^{-1}]{\mathcal{F}}\quad
    \ju\Im{X(\ju\omega)}
\shortintertext{Parseval's Relation for Aperiodic Signals}
    \int_{-\infty}^{\infty} \abs{x(t)}^2 \dd{t}
    \quad&=\quad
    \frac{1}{2\pi}
    \int_{-\infty}^{\infty} \abs{X(\ju\omega)}^2 \dd{\omega}
\end{align*}

\subsection{Symmetry}

\begin{itemize}
    \item $x(t)$ real $\quad\iff\quad$ $X(\ju\omega)=X^\ast(-\ju\omega)$
    \item $x(t)$ purely imaginary $\quad\iff\quad$ $X(\ju\omega)=-X^\ast(-\ju\omega)$
\end{itemize}

\paragraph{Conjugate Symmetry}
For real $x(t)$, its FT $X(\ju\omega)$ satisfies
\begin{itemize}
    \item $X(\ju\omega) = X^\ast(-\ju\omega)$
    \item $\Re{X(\ju\omega)} = \Re{X(-\ju\omega)}$
    \item $\Im{X(\ju\omega)} = -\Im{X(-\ju\omega)}$
    \item $\abs{X(\ju\omega)} = \abs{X(-\ju\omega)}$
    \item $\arg X(\ju\omega) = -\arg X(-\ju\omega)$
\end{itemize}

\paragraph{Parity}
\begin{itemize}
    \item $x(t)$ even $\quad\iff\quad$ $X(\ju\omega)$ even
    \item $x(t)$ odd $\quad\iff\quad$ $X(\ju\omega)$ odd
    \item $x(t)$ real \& even $\quad\implies\quad$
        $X(\ju\omega)$ real \& even
    \item $x(t)$ real \& odd $\quad\implies\quad$
        $X(\ju\omega)$ purely imaginary \& odd
\end{itemize}

\subsection{Duality}
\begin{equation*}
    f(t)
    \quad\xrightleftharpoons[\mathcal{F}^{-1}]{\mathcal{F}}\quad
    g(\omega)
    \qquad\iff\qquad
    g(t)
    \quad\xrightleftharpoons[\mathcal{F}^{-1}]{\mathcal{F}}\quad
    2\pi\cdot f(-\omega)
\end{equation*}


\section{Basic Fourier Transform Pairs}

\begin{equation*}
    \sinc(\theta) := \frac{\sin(\pi\theta)}{\pi\theta}
    \qquad\qquad
    \rect(t) := u(t+\frac{1}{2}) - u(t-\frac{1}{2})
    = \mycases{
        1 \qfor \abs{t} < \frac{1}{2} \\
        0 \qotherwise
    }
\end{equation*}

\begin{center}
    % \renewcommand{\arraystretch}{2}
\begin{tabular}
    { c @{\qquad} c @{\qquad} c}
    Signal & Fourier Transform & Fourier Series Coefficients \\
    \hline $\displaystyle
    \sum_{k=-\infty}^{\infty}
    a_k\e^{\ju k\omega_0 t}
    $ & $\displaystyle
    2\pi \sum_{k=-\infty}^{\infty} 
    a_k \delta(\omega-k\omega_0)
    $ & $
    a_k  \vphantom{\mycases{1\\1}}
    $ \\\hline $
    \e^{\ju\omega_0 t}
    $ & $\displaystyle
    2\pi \delta(\omega-\omega_0)
    $ & $
    \mycases{
        a_1 = 1 \\ a_k = 0 \qotherwise
    }
    $ \\\hline $
    \cos(\omega_0 t)
    $ & $
    \pi[ \delta(\omega-\omega_0) + \delta(\omega+\omega_0) ]
    $ & $
    \mycases{
        a_1 = a_{-1} = \frac{1}{2}\\ a_k = 0 \qotherwise
    }
    $ \\\hline $
    \sin(\omega_0 t)
    $ & $\displaystyle
    \frac{\pi}{\ju}[ \delta(\omega-\omega_0) - \delta(\omega+\omega_0) ]
    $ & $
    \mycases{
        a_1 = a_{-1} = \frac{1}{2\ju}\\ a_k = 0 \qotherwise
    }
    $ \\\hline $
    \begin{matrix}
        \text{Periodic Square Wave} \\
        \mycases{
            1 & \abs{t} < T_1 \\
            0 & T_1 < \abs{t} \leqslant T/2 \\
        } \\
        \text{ with period of } T
    \end{matrix}
    $ & $\displaystyle
    \sum_{k=-\infty}^{\infty} \frac{2\sin(k\omega_0 T_1)}{k}\delta(\omega-k\omega_0)
    $ & $\displaystyle
    \frac{\omega_0 T_1}{\pi}\sinc\qty(\frac{k\omega_0 T_1}{\pi})
    = \frac{\sin(k\omega_0 T_1)}{k\pi}
    $ \\\hline $\displaystyle
    \sum_{n=-\infty}^{+\infty}\delta(t-n T)
    $ & $\displaystyle
    \frac{2\pi}{T}\sum_{k=-\infty}^{+\infty}\delta\qty(\omega-\frac{2\pi k}{T})
    $ & $\displaystyle
    a_k = \frac{1}{T}
    $ \\\hline $\displaystyle
    \rect(\frac{t}{2 T_1}) =
    \mycases{
        1 & \abs{t} < T_1 \\
        0 & T_1 < \abs{t}
    }
    $ & $\displaystyle
    \frac{2\sin\omega T_1}{\omega}
    = 2 T_1\sinc\qty(\frac{T_1\omega}{\pi})
    $ & \\\hline $\displaystyle
    \frac{\sin W t}{\pi t}
    $ & $
    \mycases{
        1 & \abs{\omega} < W \\
        0 & W < \abs{\omega}
    }
    $ & \\\hline $
    \rect(t)
    $ & $\displaystyle
    \frac{\sin(\omega/2)}{\omega/2} = \sinc\qty(\frac{\omega}{2\pi})
    $ & \vphantom{$\mycases{1\\1}$} \\\hline $\displaystyle
    \sinc\qty(\frac{\omega}{2\pi})
    $ & $
    2\pi\rect(t)
    $ & \vphantom{$\mycases{1\\1}$} \\\hline $
    1
    $ & $
    2\pi\delta(\omega)
    $ & $
    \mycases{
        a_0 = 1 \\ a_k = 0 \qotherwise\!\!\!\!\!\!\!\!\!\!\!\!
    }
    \begin{pmatrix}
        \text{For any choice} \\
        \text{of } \omega_0 > 0
    \end{pmatrix}
    $ \\\hline $
    \delta(t)
    $ & $
    1
    $ & \vphantom{$\mycases{1\\1}$} \\\hline $
    u(t)
    $ & $\displaystyle
    \frac{1}{\ju\omega} + \pi\delta(\omega)
    $ & \vphantom{$\mycases{1\\1}$} \\\hline $
    \delta(t-t_0)
    $ & $
    \e^{-\ju\omega t_0}
    $ & \vphantom{$\mycases{1\\1}$} \\\hline $
    \e^{-a t}u(t) \qq{with} \Re{a}>0
    $ & $\displaystyle
    \frac{1}{a+\ju\omega}
    $ & \vphantom{$\mycases{1\\1}$} \\\hline $
    t\e^{-a t}u(t) \qq{with} \Re{a}>0
    $ & $\displaystyle
    \frac{1}{(a+\ju\omega)^2}
    $ & \vphantom{$\mycases{1\\1}$} \\\hline $\displaystyle
    \frac{t^{n-1}}{(n-1)!}\e^{-a t}u(t) \qq{with} \Re{a}>0
    $ & $\displaystyle
    \frac{1}{(a+\ju\omega)^n}
    $ & \vphantom{$\mycases{1\\1}$} \\\hline
\end{tabular}
\end{center}



\section{Systems Characterized By LCCDEs}

An LTI system with initial rest characterized by the LCCDE
\begin{equation*}
    \sum_{k=0}^{N} a_k \dv[k]{y(t)}{t} =
    \sum_{k=0}^{M} b_k \dv[k]{x(t)}{t}
\end{equation*}
can be represented in the frequency domain as
\begin{equation*}
    H(\ju\omega) =
    \frac{Y(\ju\omega)}{X(\ju\omega)} =
    \frac{\displaystyle\sum_{k=0}^{M} b_k (\ju\omega)^k}
    {\displaystyle\sum_{k=0}^{N} a_k (\ju\omega)^k}
\end{equation*}


\end{document}